% GENERATED_BY: oc_core_monograph_translation_draft_pipeline_v1.py
% AI_ASSISTED_TRANSLATION_DRAFT: TRUE
% THEOREM_REVIEW_STATUS: PENDING
% THEOREM_REVIEW_MODE: FORMAL_AI_GATE__CODEX_CLI_CHATGPT
% THEOREM_REVIEW_REVIEWER_ID:
% THEOREM_REVIEW_AUTHORITY_CLASS:
% THEOREM_REVIEWED_AT:
% THEOREM_REVIEW_DOSSIER_REF:
% SOURCE_LANGUAGE: EN
% TARGET_LANGUAGE: RU
% SOURCE_EN_REF: content/m_spaces/m1.tex
% ================================================================
% ==== FILE: content/m_spaces/m1.tex
% ================================================================

% ==============================
%  Ontology of Continua — Core
%  Meta-Space: M1
%  FULL MODULE — FINAL
% ==============================

\subsubsection{\texorpdfstring{$M_1$}{M_1} Обзор}
\label{sec:m1-overview}

$M_1$ — это мета-пространство, задающее структуру допустимости для
одномерного континуума $K_1$.
В отличие от $M_0$, который управляет только прото-различиями, в $M_1$
появляются первые подлинно геометрические и динамические условия допустимости:
\[
\Omega(K_1) \subseteq \Omega(M_1).
\]

$M_1$ — минимальное мета-пространство, в котором определены:

\begin{itemize}
    \item измеримая ось $A_1$;
    \item непрерывное поле $\phi(x)$ на одномерной области;
    \item прото-динамика через потоки $J_1=(\partial_x \phi, \partial_t \phi)$;
    \item допустимые энергетические функционалы и действия;
    \item корректная структурная натяжка $T_1$ и осмысленный порог $\Theta_1$;
    \item существование и правила коллапса $K_1$.
\end{itemize}

Таким образом, $M_1$ — мета-структура, делающая континуум $K_1$
математически и физически возможным.

\subsubsection*{1. Структура \texorpdfstring{$\Omega(M_1)$}{\Omega(M_1)}}

\[
\Omega(M_1)
=
\Big\{
\phi : X \to V \ \big|
\phi \in C^{0}(X,V),\
\partial_x \phi,\ \partial_t \phi\ \text{имеются в } H^1
\Big\}.
\]

$M_1$ задаёт:

\begin{enumerate}
    \item \textbf{Непрерывность:}
    Поля должны иметь непрерывные представления. Разрывные отображения
    лежат вне $\partial\Omega(M_1)$.

    \item \textbf{Соболевская регулярность:}
    Первые производные должны существовать в слабом смысле:
    \[
    \partial_x \phi,\ \partial_t \phi \in H^1.
    \]

    \item \textbf{Совместимость метрики:}
    Допустимая метрическая структура $X$ и $V$ фиксируется $M_1$ и требуется
    для вычисления энергии и натяжки.

    \item \textbf{Топологическая допустимость:}
    $X$ должен быть одномерным топологическим многообразием
    (открытый интервал, окружность, полуось и т.д.).
\end{enumerate}

Любое нарушение этих ограничений помещает конфигурацию на границу
$\partial\Omega(M_1)$, делая существование $K_1$ невозможным.

\subsubsection*{2. Допустимые оси в \texorpdfstring{$M_1$}{M_1}}

$M_1$ допускает ровно одну структурную ось:
\[
A_1 : X \to \mathbb{R},
\]
представляющую первую измеримую размерность.

В $M_1$ запрещено:

\begin{itemize}
    \item наличие дополнительных независимых осей;
    \item несовместимая нелинейная структура для $A_1$;
    \item осциллирующие оси без допустимой динамики.
\end{itemize}

Эти запреты определяют область возможных эволюций $K_1$.

\subsubsection*{3. Допустимые потенциалы и энергии}

$M_1$ требует существования корректно определённого энергетического функционала,
согласованного с универсальной операторной структурой:
\[
E[\phi] = \int_X \left(
\frac{1}{2} |\partial_x \phi|^2 + V(\phi)
\right)\,dx,
\]
где:

\begin{itemize}
    \item $V(\phi)$ ограничена снизу,
    \item коэрцитивность предотвращает коллапс,
    \item дифференцируемость для определения потоков.
\end{itemize}

Допустимые потенциалы $P(M_1)$ — это ровно те, что удовлетворяют указанным
условиям.

\subsubsection*{4. Пороги и структурная натяжка}

$M_1$ задаёт порог:
\[
\Theta_1 = \inf \{ T_1(\phi) : \phi \in \Omega(M_1) \},
\]
и функционал структурной натяжки:
\[
T_1(\phi) =
\int_X \left(
|\partial_x \phi|^2 + W(\phi)
\right)\, dx,
\]
где $W$ включает члены, обеспечивающие ограничения.

Нарушение условия $T_1 < \Theta_1$ приводит к коллапсу $K_1$.

\subsubsection*{5. Допустимые потоки и прото-динамика}

Допустимые потоки:
\[
J_1 = (\partial_x \phi,\ \partial_t \phi),
\]
с $\partial_t\phi$, принадлежащим допустимому касательному пространству
$\Omega(M_1)$.

$M_1$ запрещает потоки, которые:
\begin{itemize}
    \item увеличивают натяжку без допустимой компенсации,
    \item вызывают не-физические разрывы,
    \item нарушают принцип действия или ограничения сохранения.
\end{itemize}

\subsubsection*{6. Коллапс и граница \texorpdfstring{$M_1$}{M_1}}

Конфигурация лежит на границе $\partial\Omega(M_1)$, если:

\begin{itemize}
    \item нарушена непрерывность,
    \item нарушена Соболева регулярность,
    \item энергия становится неограниченной,
    \item $T_1 \ge \Theta_1$,
    \item $A_1$ становится неизмеримой.
\end{itemize}

Коллапс $K_1$ происходит, когда его допустимые конфигурации пересекают эту границу.

\subsubsection*{7. Связь с \texorpdfstring{$K_0$}{K_0} и \texorpdfstring{$K_2$}{K_2}}

$M_1$ — минимальное пространство, разрешающее переход $K_0 \to K_1$ через оператор
$\Psi_{0\to1}$.

Переход $K_1 \to K_2$ допустим только если:
\[
A_2 \in A(M_2),\qquad
\Omega(K_2) \subseteq \Omega(M_2),\qquad
T_1 \ge \Theta_{\text{dim}}.
\]

Следовательно, $M_1$ управляет:

\begin{itemize}
    \item существованием и устойчивостью $K_1$,
    \item допустимой областью первого динамического континуума,
    \item структурной возможностью континуумов более высокой размерности.
\end{itemize}
